% Generated by paper/master/generate_sources.py.  Do not edit.
\noindent 63 of the registered derivation statements are not recorded as proven. Each is listed with the document it lives in, its locator, and the corpus's own sentence on what remains. These are the statements an agent can act on directly.\par\medskip
\begingroup\sloppy
\noindent\textbf{\texttt{\small DERIV:\allowbreak{}EARLIER\_\allowbreak{}PROOF\_\allowbreak{}RECOVERY:\allowbreak{}DIPOLAR\_\allowbreak{}REMAINDER}} \hfill {\footnotesize open}\par
{\small Establish the archived O(|x|\textasciicircum{}-4) remainder of the dipolar lattice kernel with the required Fourier-cutoff and derivative bounds.}\par
{\footnotesize \srcpath{docs/derivations/earlier-proof-recovery.md} --- \#\# E18. Explicit remaining application obligations}\par
{\footnotesize \textit{Hypotheses.} The exact source assumptions and required estimates described in E18.}\par
{\footnotesize \textit{Remaining.} Explicit pending application or extension; neither the exact finite controls nor the narrower recovered theorem discharge it.}\par
\smallskip
\noindent\textbf{\texttt{\small DERIV:\allowbreak{}EARLIER\_\allowbreak{}PROOF\_\allowbreak{}RECOVERY:\allowbreak{}PHYSICAL\_\allowbreak{}CARRIER\_\allowbreak{}REALIZATION}} \hfill {\footnotesize open}\par
{\small Construct actual cutoff sources, residue floor, collapsing isolated intervals, true-sector gap, physical clock and limiting identifications needed to apply the positive matrix atom theorem.}\par
{\footnotesize \srcpath{docs/derivations/earlier-proof-recovery.md} --- \#\# E18. Explicit remaining application obligations}\par
{\footnotesize \textit{Hypotheses.} The exact source assumptions and required estimates described in E18.}\par
{\footnotesize \textit{Remaining.} Explicit pending application or extension; neither the exact finite controls nor the narrower recovered theorem discharge it.}\par
\smallskip
\noindent\textbf{\texttt{\small DERIV:\allowbreak{}EARLIER\_\allowbreak{}PROOF\_\allowbreak{}RECOVERY:\allowbreak{}STAPLE\_\allowbreak{}TRANSVERSALITY\_\allowbreak{}REPAIR}} \hfill {\footnotesize open}\par
{\small Repair the force-rank and tube estimates with the full moving-adjoint derivative.}\par
{\footnotesize \srcpath{docs/derivations/earlier-proof-recovery.md} --- \#\# E18. Explicit remaining application obligations}\par
{\footnotesize \textit{Hypotheses.} The exact source assumptions and required estimates described in E18.}\par
{\footnotesize \textit{Remaining.} Explicit pending application or extension; neither the exact finite controls nor the narrower recovered theorem discharge it.}\par
\smallskip
\noindent\textbf{\texttt{\small DERIV:\allowbreak{}EARLIER\_\allowbreak{}PROOF\_\allowbreak{}RECOVERY:\allowbreak{}VSU\_\allowbreak{}WHOLE\_\allowbreak{}SPACE}} \hfill {\footnotesize open}\par
{\small Establish compatible normalization and uniform local estimates for the whole-space VSU limit, then justify the claimed far-field law.}\par
{\footnotesize \srcpath{docs/derivations/earlier-proof-recovery.md} --- \#\# E18. Explicit remaining application obligations}\par
{\footnotesize \textit{Hypotheses.} The exact source assumptions and required estimates described in E18.}\par
{\footnotesize \textit{Remaining.} Explicit pending application or extension; neither the exact finite controls nor the narrower recovered theorem discharge it.}\par
\smallskip
\noindent\textbf{\texttt{\small DERIV:\allowbreak{}UNIVERSAL\_\allowbreak{}CELLULAR\_\allowbreak{}HODGE:\allowbreak{}PROPER\_\allowbreak{}RETURN\_\allowbreak{}Q\_\allowbreak{}ANNIHILATION}} \hfill {\footnotesize open}\par
{\small The 60 paths flagged by the submitted endpoint-or-zero flux predicate identify physical tetrahedral retained-sector returns that vanish individually under intermediate Q projection.}\par
{\footnotesize \srcpath{docs/derivations/universal-cellular-hodge-tetrahedral.md} --- \#\#\# 3.4 Open physical-history identification (lines 133-156)}\par
{\footnotesize \textit{Hypotheses.} An actual Haar/Fierz Hilbert space, physical retained projector, insertion channels, and projected resolvent states matching the enumerated paths.}\par
{\footnotesize \textit{Remaining.} Open: construct and identify the physical states/projector and calculate each projected contribution. The 60/36 flux count and Q psi=0 do not establish this identification.}\par
\smallskip
\noindent\textbf{\texttt{\small DERIV:\allowbreak{}W6\_\allowbreak{}CONDITIONAL\_\allowbreak{}SCORE\_\allowbreak{}TAIL\_\allowbreak{}CONTROL:\allowbreak{}SCORE\_\allowbreak{}DOMINATION\_\allowbreak{}M10}} \hfill {\footnotesize open}\par
{\small For the actual compact square and specified source-preserving dilation, prove K\_\allowbreak{}g(w)<=C0+C1 W\_\allowbreak{}g(w) almost everywhere in nu\_\allowbreak{}g, where W\_\allowbreak{}g(w)=g\textasciicircum{}-2 E\_\allowbreak{}(mu\_\allowbreak{}g)(V|w), with finite nonnegative C0,C1 independent of 0<g<g\_\allowbreak{}*. This sufficient conditional-score domination is open.}\par
{\footnotesize \srcpath{docs/derivations/w6-conditional-score-tail-control.md} --- \#\#\# M.4. The precise remaining conditional-score inequality (lines 473-495)}\par
{\footnotesize \textit{Hypotheses.} Actual fixed twelve-edge, four-face compact Wilson square, true normalized quantum ground Psi\_\allowbreak{}g, dmu\_\allowbreak{}g=Psi\_\allowbreak{}g\textasciicircum{}2 dU, literal source w=Tr(U2 U3)/2 and its marginal nu\_\allowbreak{}g.; sigma\_\allowbreak{}g=(partial\_\allowbreak{}g Psi\_\allowbreak{}g+D Psi\_\allowbreak{}g/g)/Psi\_\allowbreak{}g uses the specified compact dilation of Q8, including its Haar divergence; K\_\allowbreak{}g=Var\_\allowbreak{}(mu\_\allowbreak{}g)(sigma\_\allowbreak{}g|w) retains all compact corrections and rare coarse fibers.; The established actual fixed-square gap and ground-energy bounds hold on the specified small-coupling interval. Averaged conditional-score control is not assumed to imply pointwise domination.}\par
{\footnotesize \textit{Remaining.} Prove M10 for the actual conditional law, or retain its open status if a different sufficient source-energy estimate is proved. The phase-adapted local implication requires whole-fiber tails, conditional-mean gluing and consistent global dilation before it can address this target.}\par
\smallskip
\noindent\textbf{\texttt{\small DERIV:\allowbreak{}W6\_\allowbreak{}CONDITIONAL\_\allowbreak{}TRANSPORT\_\allowbreak{}OBSTRUCTION:\allowbreak{}M10\_\allowbreak{}SYNCHRONIZED\_\allowbreak{}SCORE}} \hfill {\footnotesize open}\par
{\small The selected synchronized realization of the generic M10 target asks whether the actual score for the specified S13 compact field obeys K\_\allowbreak{}g(w)<=C0+C1 g\textasciicircum{}-2 E\_\allowbreak{}mu\_\allowbreak{}g[V|w] with finite nonnegative constants uniform in sufficiently small positive g and almost every retained trace w.}\par
{\footnotesize \srcpath{docs/derivations/w6-conditional-transport-obstruction.md} --- \#\# S6 Remaining all-fiber estimate (lines 317-375), explicitly open synchronized M10 target}\par
{\footnotesize \textit{Hypotheses.} Actual positive four-face SU(2) square ground and literal retained trace w=tr(U2 U3)/2.; The particular synchronized four radial profiles S13 and their full Haar unitary generator, not an arbitrary Q8 field.; All conditional source regimes, including rare fibers and the antipodal degeneration, are retained in the requested uniform estimate.}\par
{\footnotesize \textit{Remaining.} Prove the actual full conditional score bound, or a sufficient direct source-weighted pairing estimate. The original M11-M15 consequence then becomes applicable;}\par
\smallskip
\noindent\textbf{\texttt{\small DERIV:\allowbreak{}W6\_\allowbreak{}GROUND\_\allowbreak{}JETS\_\allowbreak{}TRANSPORT\_\allowbreak{}BUDGET:\allowbreak{}INTERACTING\_\allowbreak{}GRID\_\allowbreak{}COMPARISON}} \hfill {\footnotesize open}\par
{\small For the actual interacting growing-grid family, obtain a volume-independent fast floor and complete source/transport energy estimates using linked or locally centered energy bounds, and control the coarse-theory matching discrepancy after the prescribed injection. The fixed-square constants depending on total ground energy E and the Gaussian grid floor do not prove this extension.}\par
{\footnotesize \srcpath{docs/derivations/w6-ground-jets-and-transport-budget.md} --- \#\# 6. Source energies, subdivision, interacting grids, and the physical clock (lines 351-366)}\par
{\footnotesize \textit{Hypotheses.} Actual coupled Wilson grids and specified block/source injections, with true vacuum subtraction, common physical energy normalization and the required retained, fast and spectral-energy domains.; Constants must be independent of growing spatial volume in the intended small-coupling regime; the existing fixed-square E and gamma are not assumed to have that uniformity.; Actual source-energy comparison factors and the discrepancy between the generated coarse form and prescribed coarse theory are retained; the Gaussian covariance-weighted source is not identified with the nonlinear quantum source without proof.}\par
{\footnotesize \textit{Remaining.} Prove the interacting fast-floor and connected energy/source comparison in the specified grid family. Controlling the local coupling remainder does not control coarse matching, the explicit physical-clock change in R14, or convergence of the exact transported jets and observables.}\par
\smallskip
\noindent\textbf{\texttt{\small DERIV:\allowbreak{}W6\_\allowbreak{}GROUND\_\allowbreak{}JETS\_\allowbreak{}TRANSPORT\_\allowbreak{}BUDGET:\allowbreak{}SOURCE\_\allowbreak{}ENERGY\_\allowbreak{}JETS\_\allowbreak{}R10}} \hfill {\footnotesize open}\par
{\small For the actual complete compact-square source/vacuum generator A(g) in R9, prove ||A\textasciicircum{}(r)(g)||\_\allowbreak{}(q\_\allowbreak{}g->q\_\allowbreak{}g)<=d\_\allowbreak{}r g\textasciicircum{}(-r-1) for r=0,1,2 with finite constants independent of small positive g. This is the open model input R10, not the already proved conditional implication R10 to R12.}\par
{\footnotesize \srcpath{docs/derivations/w6-ground-jets-and-transport-budget.md} --- \#\# 4. An explicit bridge from the actual transport generator to M1,M2,M3 (lines 199-218); section 5 (lines 271-309)}\par
{\footnotesize \textit{Hypotheses.} Actual fixed compact-square Hamiltonian and common electric H1 form domain; actual simple normalized positive ground, 0<=e\_\allowbreak{}g<=E and physical gap gamma>0 on 0<g<g\_\allowbreak{}*.; P\_\allowbreak{}g is the complete true-ground literal-source projection from C1, and A(g) is exactly the R9 source/vacuum generator; q\_\allowbreak{}g[u]=H\_\allowbreak{}g[u]+gamma||u||\textasciicircum{}2, including vacuum components.; Parameter derivatives are taken at positive g on the fixed compact form domain. The target covers all source and fast energies, including spatial derivatives and generator derivatives through order two.}\par
{\footnotesize \textit{Remaining.} Establish the conditional-source projection energy bounds beyond the proved vacuum-only generator. A direct proof of mixed-form R11 is an alternative sufficient route and need not prove R10 itself.}\par
\smallskip
\noindent\textbf{\texttt{\small DERIV:\allowbreak{}W6\_\allowbreak{}SOURCE\_\allowbreak{}ENERGY\_\allowbreak{}JETS\_\allowbreak{}R10:\allowbreak{}H0\_\allowbreak{}ENERGY\_\allowbreak{}CONTRACTION}} \hfill {\footnotesize open}\par
{\small Establish the proposed bound b\_\allowbreak{}g[Pi\_\allowbreak{}g Phi] <= q\_\allowbreak{}g[Omega\_\allowbreak{}g Phi] with kappa\_\allowbreak{}P=1 on the actual model; the weaker uniform finite-kappa\_\allowbreak{}P H0 also suffices for R10.}\par
{\footnotesize \srcpath{docs/derivations/w6-source-energy-jets-r10.md} --- \#\# 6. Scaling defects and unresolved estimates (C6)}\par
{\footnotesize \textit{Hypotheses.} Actual fixed compact square at positive g; normalized positive ground, q\_\allowbreak{}g=H\_\allowbreak{}g+gamma I, common form domain.; Constants must be independent of g on 0<g<g\_\allowbreak{}*.}\par
{\footnotesize \textit{Remaining.} The submitted proof is not a discharge. C1-C5 repair identities and conditional reductions;}\par
\smallskip
\noindent\textbf{\texttt{\small DERIV:\allowbreak{}W6\_\allowbreak{}SOURCE\_\allowbreak{}ENERGY\_\allowbreak{}JETS\_\allowbreak{}R10:\allowbreak{}H1\_\allowbreak{}FIBERWISE\_\allowbreak{}SCORE\_\allowbreak{}VARIANCE}} \hfill {\footnotesize open}\par
{\small Establish sup\_\allowbreak{}w Var\_\allowbreak{}mu\_\allowbreak{}g(sigma\_\allowbreak{}g|w) <= kappa\_\allowbreak{}0 g\textasciicircum{}-2 uniformly at small positive coupling.}\par
{\footnotesize \srcpath{docs/derivations/w6-source-energy-jets-r10.md} --- \#\# 6. Scaling defects and unresolved estimates (C6)}\par
{\footnotesize \textit{Hypotheses.} Actual fixed compact square at positive g; normalized positive ground, q\_\allowbreak{}g=H\_\allowbreak{}g+gamma I, common form domain.; Constants must be independent of g on 0<g<g\_\allowbreak{}*.}\par
{\footnotesize \textit{Remaining.} The submitted proof is not a discharge. C1-C5 repair identities and conditional reductions;}\par
\smallskip
\noindent\textbf{\texttt{\small DERIV:\allowbreak{}W6\_\allowbreak{}SOURCE\_\allowbreak{}ENERGY\_\allowbreak{}JETS\_\allowbreak{}R10:\allowbreak{}KATO\_\allowbreak{}DERIVATIVES}} \hfill {\footnotesize open}\par
{\small Establish ||partial\_\allowbreak{}g\textasciicircum{}r[P\_\allowbreak{}g prime,P\_\allowbreak{}g]||\_\allowbreak{}(q\_\allowbreak{}g->q\_\allowbreak{}g) <= d\_\allowbreak{}P,r g\textasciicircum{}(-r-1), r=0,1,2, for the actual source projection.}\par
{\footnotesize \srcpath{docs/derivations/w6-source-energy-jets-r10.md} --- \#\# 6. Scaling defects and unresolved estimates (C6)}\par
{\footnotesize \textit{Hypotheses.} Actual fixed compact square at positive g; normalized positive ground, q\_\allowbreak{}g=H\_\allowbreak{}g+gamma I, common form domain.; Constants must be independent of g on 0<g<g\_\allowbreak{}*.}\par
{\footnotesize \textit{Remaining.} The submitted proof is not a discharge. C1-C5 repair identities and conditional reductions;}\par
\smallskip
\noindent\textbf{\texttt{\small DERIV:\allowbreak{}W6\_\allowbreak{}SOURCE\_\allowbreak{}ENERGY\_\allowbreak{}JETS\_\allowbreak{}R10:\allowbreak{}VACUUM\_\allowbreak{}CROSS\_\allowbreak{}DERIVATIVES}} \hfill {\footnotesize open}\par
{\small Establish ||partial\_\allowbreak{}g\textasciicircum{}r A\_\allowbreak{}nu||\_\allowbreak{}(q\_\allowbreak{}g->q\_\allowbreak{}g) <= d\_\allowbreak{}nu,r g\textasciicircum{}(-r-1), r=0,1,2, for the actual vacuum-cross source vector.}\par
{\footnotesize \srcpath{docs/derivations/w6-source-energy-jets-r10.md} --- \#\# 6. Scaling defects and unresolved estimates (C6)}\par
{\footnotesize \textit{Hypotheses.} Actual fixed compact square at positive g; normalized positive ground, q\_\allowbreak{}g=H\_\allowbreak{}g+gamma I, common form domain.; Constants must be independent of g on 0<g<g\_\allowbreak{}*.}\par
{\footnotesize \textit{Remaining.} The submitted proof is not a discharge. C1-C5 repair identities and conditional reductions;}\par
\smallskip
\noindent\textbf{\texttt{\small DERIV:\allowbreak{}W6\_\allowbreak{}SOURCE\_\allowbreak{}GENERATOR\_\allowbreak{}SCORE\_\allowbreak{}FRAME:\allowbreak{}CONDITIONAL\_\allowbreak{}ENERGY\_\allowbreak{}H0\_\allowbreak{}H4}} \hfill {\footnotesize open}\par
{\small Prove, uniformly on 0<g<g\_\allowbreak{}*, the fiberwise conditional-energy estimates H0-H4 for the undilated ground score on level sets of w. Their averaged forms are established (SF7b and the integrated J\_\allowbreak{}g bound conditional on H4); the fiberwise forms are open.}\par
{\footnotesize \srcpath{docs/derivations/w6-source-generator-score-frame.md} --- \#\# 5. Exact reduction of the order-zero R10 bound (SF8-SF9), conditional-energy hypotheses}\par
{\footnotesize \textit{Hypotheses.} Actual fixed square, undilated score partial\_\allowbreak{}g log Omega\_\allowbreak{}g; K\textasciicircum{}0\_\allowbreak{}g is not the dilated K\_\allowbreak{}g of M10.}\par
{\footnotesize \textit{Remaining.} A fiberwise Poincare-type inequality for -Delta\_\allowbreak{}mu restricted to level sets of w, uniform in g, would give H1 and H2 from SF6. H0, H3 and H4 concern the source energy of conditional expectations and need a separate argument.}\par
\smallskip
\noindent\textbf{\texttt{\small DERIV:\allowbreak{}W6\_\allowbreak{}SYNCHRONIZED\_\allowbreak{}M10\_\allowbreak{}DOMINATION:\allowbreak{}AMPLITUDE\_\allowbreak{}GRADIENT\_\allowbreak{}SCALING}} \hfill {\footnotesize open}\par
{\small The actual uniform amplitude-score gradient bound ca/g remains open; changing parameter to h=g\textasciicircum{}2 and a spatial WKB expansion do not alone prove it.}\par
{\footnotesize \srcpath{docs/derivations/w6-synchronized-m10-domination.md} --- \#\# 5. Correct parameter-derivative criterion (lines 144-171)}\par
{\footnotesize \textit{Hypotheses.} Actual positive ground amplitude in a valid phase tube.}\par
{\footnotesize \textit{Remaining.} Establish uniform spatial control of h*partial\_\allowbreak{}h log A and Z log A for the actual ground; the logarithmic-parameter criterion gives the sufficient implication.}\par
\smallskip
\noindent\textbf{\texttt{\small DERIV:\allowbreak{}W6\_\allowbreak{}SYNCHRONIZED\_\allowbreak{}M10\_\allowbreak{}DOMINATION:\allowbreak{}FULL\_\allowbreak{}SCORE\_\allowbreak{}DOMINATION\_\allowbreak{}M10}} \hfill {\footnotesize open}\par
{\small Actual full-fiber M10 for S13 remains open. The earlier completion claim is retracted; the repaired assembly implication retains all unsupplied model hypotheses.}\par
{\footnotesize \srcpath{docs/derivations/w6-synchronized-m10-domination.md} --- \#\# 7. Assembly criterion and actual status (lines 262-286)}\par
{\footnotesize \textit{Hypotheses.} Actual four-face compact SU(2) quantum ground law and S13 dilation.}\par
{\footnotesize \textit{Remaining.} Actual complement is discharged. Supply the polynomial relative-amplitude and tube moment estimates, intermediate controls, and actual normal/angular comparison and score error.}\par
\smallskip
\noindent\textbf{\texttt{\small DERIV:\allowbreak{}WILSON\_\allowbreak{}MARKED\_\allowbreak{}TRANSFER:\allowbreak{}WT6}} \hfill {\footnotesize open}\par
{\small The proposed rooted coefficient-operator bound WT6 is an explicit sufficient route to complete-shell transport; this historical argument did not establish it. Its continuation uses a different full-Hilbert estimate and finite-order locality.}\par
{\footnotesize \srcpath{docs/derivations/wilson-marked-transfer.md} --- \#\# WT-6. The remaining complete-shell estimate and conditional transport (lines 209-248)}\par
{\footnotesize \textit{Hypotheses.} Finite cubic link Hilbert product; positive free generators kill the product vacuum and have gap delta>0 on its complement.; Bounded plaquette interactions ||V\_\allowbreak{}p||<=v; physical character-clock step tau>0 and blocks of duration ell=m*tau.; Actual vacuum-dressed transfer on the specified exact-support coefficient space must first be constructed.}\par
{\footnotesize \textit{Remaining.} Formalize or discharge the particular WT6 rooted norm if used; do not equate the later sufficient route with a proof of this specific estimate.}\par
\smallskip
\noindent\textbf{\texttt{\small DERIV:\allowbreak{}WILSON\_\allowbreak{}SC17\_\allowbreak{}SPATIAL\_\allowbreak{}CLOSURE:\allowbreak{}ACTUAL\_\allowbreak{}REFERENCE\_\allowbreak{}DEFECT\_\allowbreak{}BOUND}} \hfill {\footnotesize open}\par
{\small On the specified actual Wilson conditional source chart, establish a volume-uniform weighted bound on the complete signed Hessian defect D\_\allowbreak{}B=Hess(V\_\allowbreak{}B+(H\_\allowbreak{}out Omega)/Omega)-2epsilon A\textasciicircum{}2, including the actual kinetic metric, moving source projector and localization terms, with norm(D\_\allowbreak{}B)\_\allowbreak{}(s,infty)<=delta and 2 beta\_\allowbreak{}s(A)\textasciicircum{}2 delta/epsilon<1.}\par
{\footnotesize \srcpath{docs/derivations/wilson-sc17-spatial-closure.md} --- \#\# 4. The actual quantum pressure and cutoff defects are part of the input (lines 698-756), with the R5-R7 budget in section 2}\par
{\footnotesize \textit{Hypotheses.} Actual positive Wilson vacuum, compatible source and block coordinates, and multiplication terms assigned exactly once; the complete coupled Gaussian reference is subtracted.; The reference semigroup budget beta\_\allowbreak{}s(A) is finite in the chosen weighted norm at the fixed blocking scale; epsilon>0 and all conditional, metric and cutoff domains are retained.; The desired uniformity concerns the stated block family and surrounding volume; a separate source-compatible reference-angle margin is still needed for global assembly.}\par
{\footnotesize \textit{Remaining.} Prove the actual complete weighted smallness estimate. The proved Gaussian pressure cancellation, an isolated magnetic Taylor coefficient, or formalizing the conditional Riccati implication does not discharge this model obligation.}\par
\smallskip
\noindent\textbf{\texttt{\small DERIV:\allowbreak{}WILSON\_\allowbreak{}SELECTED\_\allowbreak{}INVERSE\_\allowbreak{}WALL:\allowbreak{}W6}} \hfill {\footnotesize open}\par
{\small The needed actual interacting Wilson bound is sup\_\allowbreak{}(Lambda,b[p]=1)|d\_\allowbreak{}g,QQ[R0 t1p,Rg t1p]|<=C|g| uniformly in volume, coupling, backgrounds and required spectral interval.}\par
{\footnotesize \srcpath{docs/derivations/wilson-selected-inverse-wall.md} --- \#\# 5. The exact point at which this argument cannot proceed (lines 221-264)}\par
{\footnotesize \textit{Hypotheses.} The operator chart, vacuum subtraction and fast/retained spaces are those of the spatial Schur derivation; domains are explicitly checked rather than identified by notation.; 0<g<=g0; complete relevant retained energy space and high-retained contribution controlled.}\par
{\footnotesize \textit{Remaining.} Prove the actual coupled signed resolvent pairing with the complete dressed force, physical energy normalization and an interacting fast floor; no existing Gaussian control discharges it.}\par
\smallskip
\noindent\textbf{\texttt{\small DERIV:\allowbreak{}YANGMILLS\_\allowbreak{}CONTINUUM\_\allowbreak{}BALABAN\_\allowbreak{}MULTISCALE\_\allowbreak{}PROOF:\allowbreak{}OS1}} \hfill {\footnotesize open}\par
{\small Full Euclidean covariance requires restoration of SO(4) beyond cubic block symmetry; the repair record ties this to the unproved irrelevance estimate.}\par
{\footnotesize \srcpath{docs/derivations/yangmills-continuum-balaban-multiscale-proof.md} --- \#\#\# 8.2 OS1: Euclidean Covariance (lines 375-382)}\par
{\footnotesize \textit{Hypotheses.} Submitted SU(N), N>=2 continuum multiscale construction; read together with its in-document Repair record audited 2026-09-09.; The source title and summary claim completion, but its repair record retains explicit averaging/interpolation, irrelevance, physical-scale and scattering obligations. No machine certification is inferred from prose.}\par
{\footnotesize \textit{Remaining.} Prove rotational restoration for the actual continuum limit, then translation/rotation covariance of its Schwinger distributions.}\par
\smallskip
\noindent\textbf{\texttt{\small DERIV:\allowbreak{}YANGMILLS\_\allowbreak{}CONTINUUM\_\allowbreak{}BALABAN\_\allowbreak{}MULTISCALE\_\allowbreak{}PROOF:\allowbreak{}THEOREM\_\allowbreak{}7\_\allowbreak{}1}} \hfill {\footnotesize open}\par
{\small The submitted Cauchy theorem needs a summable cross-scale expectation increment. The printed O(1/k) bound alone is not summable; the repair record names the missing irrelevance factor a\_\allowbreak{}k\textasciicircum{}theta.}\par
{\footnotesize \srcpath{docs/derivations/yangmills-continuum-balaban-multiscale-proof.md} --- \#\#\# 7.1 Cauchy Sequence of Effective Actions (lines 317-340)}\par
{\footnotesize \textit{Hypotheses.} Submitted SU(N), N>=2 continuum multiscale construction; read together with its in-document Repair record audited 2026-09-09.; The source title and summary claim completion, but its repair record retains explicit averaging/interpolation, irrelevance, physical-scale and scattering obligations. No machine certification is inferred from prose.; Actual gauge-invariant cylinder expectations and uniform operator normalization; a positive irrelevance exponent or another summable increment bound.}\par
{\footnotesize \textit{Remaining.} Prove the actual summable remainder after marginal matching, then construct the Cauchy limit with its observable norms and scale maps.}\par
\smallskip
\noindent\textbf{\texttt{\small DERIV:\allowbreak{}YANGMILLS\_\allowbreak{}CONTINUUM\_\allowbreak{}BALABAN\_\allowbreak{}MULTISCALE\_\allowbreak{}PROOF:\allowbreak{}THEOREM\_\allowbreak{}8\_\allowbreak{}2}} \hfill {\footnotesize open}\par
{\small The claimed continuum exponential clustering requires a strictly positive rate in common physical units at a fixed physical scale; UV cell-scale polymer decay alone does not provide that conclusion.}\par
{\footnotesize \srcpath{docs/derivations/yangmills-continuum-balaban-multiscale-proof.md} --- \#\#\# 8.5 OS4: Exponential Clustering and Mass Gap (lines 412-436)}\par
{\footnotesize \textit{Hypotheses.} Submitted SU(N), N>=2 continuum multiscale construction; read together with its in-document Repair record audited 2026-09-09.; The source title and summary claim completion, but its repair record retains explicit averaging/interpolation, irrelevance, physical-scale and scattering obligations. No machine certification is inferred from prose.; The repair record explicitly requires the missing bridge to a scale with g=O(1); uniform cutoff/volume estimates are retained.}\par
{\footnotesize \textit{Remaining.} Prove the actual nonperturbative KP or equivalent correlation estimate at a fixed physical reference scale, then transport it through the constructed continuum limit.}\par
\smallskip
\noindent\textbf{\texttt{\small DERIV:\allowbreak{}YANGMILLS\_\allowbreak{}CONTINUUM\_\allowbreak{}BALABAN\_\allowbreak{}MULTISCALE\_\allowbreak{}PROOF:\allowbreak{}GAUGE\_\allowbreak{}HESSIAN}} \hfill {\footnotesize disputed}\par
{\small The submission gives the background gauge Hessian (5.3)-(5.5); the operator must include the correct covariant 1-form curvature term and the repaired averaging map consistently.}\par
{\footnotesize \srcpath{docs/derivations/yangmills-continuum-balaban-multiscale-proof.md} --- \#\#\# 5.1 Fluctuation Gauge Fixing (lines 183-197)}\par
{\footnotesize \textit{Hypotheses.} Submitted SU(N), N>=2 continuum multiscale construction; read together with its in-document Repair record audited 2026-09-09.; The source title and summary claim completion, but its repair record retains explicit averaging/interpolation, irrelevance, physical-scale and scattering obligations. No machine certification is inferred from prose.; Smooth background and gauge-fixed fluctuation domain; repaired deterministic block constraint.}\par
{\footnotesize \textit{Remaining.} Formalize the covariant exterior-derivative/Weitzenbock identity with signs and domains; reconcile Q\_\allowbreak{}k versus Q\_\allowbreak{}B before using the Hessian.}\par
\smallskip
\noindent\textbf{\texttt{\small DERIV:\allowbreak{}YANGMILLS\_\allowbreak{}CONTINUUM\_\allowbreak{}BALABAN\_\allowbreak{}MULTISCALE\_\allowbreak{}PROOF:\allowbreak{}PROPOSITION\_\allowbreak{}3\_\allowbreak{}4}} \hfill {\footnotesize disputed}\par
{\small Proposition 3.4 asserts a unique gauge-equivariant constrained background minimizer with the Sobolev/coercivity bounds (3.9)-(3.12); the repair record says the corrected affine adjoint constraint and transported interpolation have not been propagated.}\par
{\footnotesize \srcpath{docs/derivations/yangmills-continuum-balaban-multiscale-proof.md} --- \#\#\# 3.2 The Background Field Variational Problem (lines 111-139)}\par
{\footnotesize \textit{Hypotheses.} Submitted SU(N), N>=2 continuum multiscale construction; read together with its in-document Repair record audited 2026-09-09.; The source title and summary claim completion, but its repair record retains explicit averaging/interpolation, irrelevance, physical-scale and scattering obligations. No machine certification is inferred from prose.; Small curvature g\_\allowbreak{}(k+1)\textasciicircum{}kappa a\_\allowbreak{}(k+1)\textasciicircum{}-2, 0<kappa<1/6, sufficiently small coupling; all correct gauge and boundary domains still required.}\par
{\footnotesize \textit{Remaining.} First implement the actual repaired deterministic constraint and equivariant interpolation, then prove coercivity, existence, uniqueness and smooth parameter dependence on its genuine domain.}\par
\smallskip
\noindent\textbf{\texttt{\small DERIV:\allowbreak{}YANGMILLS\_\allowbreak{}CONTINUUM\_\allowbreak{}BALABAN\_\allowbreak{}MULTISCALE\_\allowbreak{}PROOF:\allowbreak{}THEOREM\_\allowbreak{}5\_\allowbreak{}1}} \hfill {\footnotesize disputed}\par
{\small Theorem 5.1 asserts a uniform positive background fluctuation floor and exponential pointwise resolvent kernel bound. The actual averaging conjugation, covariant Poincare and kernel regularity estimates remain required.}\par
{\footnotesize \srcpath{docs/derivations/yangmills-continuum-balaban-multiscale-proof.md} --- \#\#\# 5.3 Combes--Thomas Resolvent Localization for Continuous \$\textbackslash{}mathcal\{H\}\_\allowbreak{}\{B\_\allowbreak{}k\}\$ (lines 209-248)}\par
{\footnotesize \textit{Hypotheses.} Submitted SU(N), N>=2 continuum multiscale construction; read together with its in-document Repair record audited 2026-09-09.; The source title and summary claim completion, but its repair record retains explicit averaging/interpolation, irrelevance, physical-scale and scattering obligations. No machine certification is inferred from prose.; delta<=1/(4224ell\textasciicircum{}2); alpha0>=8delta+1; the domain, cutoff and averaging operator must be specified.}\par
{\footnotesize \textit{Remaining.} Formalize the corrected covariant averaged operator and exponential conjugation including its nonlocal averaging term; replace or justify the all-x,y bounded continuum resolvent kernel claim, which has a diagonal singularity without an explicit smoothing cutoff.}\par
\smallskip
\noindent\textbf{\texttt{\small DERIV:\allowbreak{}YANGMILLS\_\allowbreak{}RESOLVENT\_\allowbreak{}LOCALIZATION\_\allowbreak{}AND\_\allowbreak{}DIRICHLET\_\allowbreak{}GAP:\allowbreak{}SUBMITTED\_\allowbreak{}W6\_\allowbreak{}CLOSURE}} \hfill {\footnotesize disputed}\par
{\small The proposal asserts actual Wilson W6 closure from massive spatial Green decay; its own source audit records that the computed massive scalar operator is not identified with Wilson.}\par
{\footnotesize \srcpath{docs/derivations/yangmills-resolvent-localization-and-dirichlet-gap.md} --- \#\# 1. Summary of Results (lines 20-39)}\par
{\footnotesize \textit{Hypotheses.} This submitted proposal is explicitly preserved under its September 9 source audit; closure claims are not the current graph verdict.; The actual model/operator identification is an obligation, separate from a scalar finite replay.}\par
{\footnotesize \textit{Remaining.} Establish the actual Wilson fast floor, connected signed pairing and physical source identification; the submission does not discharge these.}\par
\smallskip
\noindent\textbf{\texttt{\small DERIV:\allowbreak{}EARLIER\_\allowbreak{}PROOF\_\allowbreak{}RECOVERY:\allowbreak{}B6\_\allowbreak{}RETAINED\_\allowbreak{}KRYLOV\_\allowbreak{}CAMPAIGN}} \hfill {\footnotesize conditional}\par
{\small The retained SU(3) B6 certificate reports radial dimensions 67,155,133 and maximum full-space Ritz residual 8.775675697349887e-14 on 201 g-values in [1,2], with crossing g=1.3039546641713.}\par
{\footnotesize \srcpath{docs/derivations/earlier-proof-recovery.md} --- \#\# E6b. Retained B6 numerical campaign}\par
{\footnotesize \textit{Hypotheses.} Historical source K6(u)=E-uM, source coupling conversion and branch choices; level-eight 511-column word spaces; relative SVD cutoff 1e-10.; The certificate is retained and hash-verified; the B6 eigensolver was not rerun in this integration.}\par
{\footnotesize \textit{Remaining.} The original numerical certificate is preserved. A fresh B6 eigensolver replay and any uniform-in-coupling or full-shell conclusion are not supplied.}\par
\smallskip
\noindent\textbf{\texttt{\small DERIV:\allowbreak{}W6\_\allowbreak{}CONDITIONAL\_\allowbreak{}SCORE\_\allowbreak{}TAIL\_\allowbreak{}CONTROL:\allowbreak{}SCORE\_\allowbreak{}WEIGHT\_\allowbreak{}FROM\_\allowbreak{}M10}} \hfill {\footnotesize conditional}\par
{\small Conditional on the actual M10 domination, every centered finite-energy source satisfies integral K\_\allowbreak{}g|f|\textasciicircum{}2 dnu\_\allowbreak{}g <= [C1+(C0+C1 E)/gamma] b\_\allowbreak{}g[f], and the actual median Hardy constant obeys mathfrak\_\allowbreak{}B\_\allowbreak{}g<=C1+2(C0+C1 E)/gamma uniformly at small coupling. The implication is analytically proved; its M10 model input remains open.}\par
{\footnotesize \srcpath{docs/derivations/w6-conditional-score-tail-control.md} --- \#\#\# M.4. The precise remaining conditional-score inequality and M.5 (lines 497-563)}\par
{\footnotesize \textit{Hypotheses.} Actual compact-square source form b\_\allowbreak{}g[f]=g\textasciicircum{}2 integral (1-w\textasciicircum{}2)|f'|\textasciicircum{}2 dnu\_\allowbreak{}g, actual ground-energy bound e\_\allowbreak{}g<=E and physical gap gamma>0, with the true-ground source identity M7.; The actual conditional-score domination M10 holds with finite nonnegative C0,C1 uniform in g.; Sources belong to the full compact retained form domain; centered sources give M12, and the median-anchored half-source argument retains the uncentered vacuum-energy term for M14-M15.}\par
{\footnotesize \textit{Remaining.} Formalize this already established conditional implication if needed. Actual application requires M10 or another sufficient score estimate;}\par
\smallskip
\noindent\textbf{\texttt{\small DERIV:\allowbreak{}W6\_\allowbreak{}CONDITIONAL\_\allowbreak{}TRANSPORT\_\allowbreak{}OBSTRUCTION:\allowbreak{}LOCAL\_\allowbreak{}GLOBAL\_\allowbreak{}SCORE\_\allowbreak{}CRITERION}} \hfill {\footnotesize conditional}\par
{\small The stated uniform conditional tube moments, differentiated amplitude-score bound and synchronized phase remainder imply the local S14 score estimate. S15 glues it to the full conditional variance using the actual outside-tube squared deviation from the same reference beta\_\allowbreak{}g. Uniform potential-weighted complement and remaining-regime bounds would supply synchronized M10.}\par
{\footnotesize \srcpath{docs/derivations/w6-conditional-transport-obstruction.md} --- \#\# S6 Remaining all-fiber estimate (lines 317-375)}\par
{\footnotesize \textit{Hypotheses.} Conditional tube moments E|eta|\textasciicircum{}(2j)<=c\_\allowbreak{}(2j)g\textasciicircum{}(2j), j=1,2,3, with uniform constants.; True relative amplitude score has fast derivative at most c\_\allowbreak{}a/g; the phase is smooth with bounded third fast derivatives and the stated oscillator quadratic jet.; Synchronized tangency holds. The same reference beta\_\allowbreak{}g is used inside and outside the tube, so conditional mean differences remain in the outside deviation.; Full M10 additionally needs a uniform potential-weighted outside deviation bound and a degeneration-uniform actual score bound in the remaining source regimes.}\par
{\footnotesize \textit{Remaining.} Establish the actual uniform differentiated amplitude, tube-moment, conditional complement and antipodal estimates. Undifferentiated local comparison used for the obstruction does not discharge these hypotheses.}\par
\smallskip
\noindent\textbf{\texttt{\small DERIV:\allowbreak{}W6\_\allowbreak{}GROUND\_\allowbreak{}JETS\_\allowbreak{}TRANSPORT\_\allowbreak{}BUDGET:\allowbreak{}COMPLETE\_\allowbreak{}BUDGET}} \hfill {\footnotesize conditional}\par
{\small For the complete actual positive-reference source/vacuum transport, energy-operator bounds ||A\textasciicircum{}(r)(g)||\_\allowbreak{}(q\_\allowbreak{}g->q\_\allowbreak{}g)<=d\_\allowbreak{}r g\textasciicircum{}(-r-1), r=0,1,2, imply explicit all-energy M\_\allowbreak{}j(s)<=c\_\allowbreak{}j s\textasciicircum{}-j for j=1,2,3 and hence the complete compact Schur residual budget. The composite budget retains changing source energies and physical clock factors.}\par
{\footnotesize \srcpath{docs/derivations/w6-ground-jets-and-transport-budget.md} --- \#\# 4. An explicit bridge from the actual transport generator to M1,M2,M3 (lines 199-367)}\par
{\footnotesize \textit{Hypotheses.} Complete actual common-domain source/vacuum transport R'=A R and nonreducing graph from the positive-reference source theorem.; For q\_\allowbreak{}g=H\_\allowbreak{}g+gamma, full energy-space ||A\textasciicircum{}(r)(g)||<=d\_\allowbreak{}r g\textasciicircum{}(-r-1), r=0,1,2, with finite uniform constants.; Local t in [s/2,s], fixed spectral window below the actual gap. The composite source growth and subdivision inputs are the prior S2-S5 theorem.; Positive physical clock tau\_\allowbreak{}n; physical matching discrepancies are separately retained, not assumed zero.}\par
{\footnotesize \textit{Remaining.} Prove the actual complete conditional source-generator energy jets or exact mixed-form substitute; supply interacting-grid, coarse-matching and physical-clock convergence inputs.}\par
\smallskip
\noindent\textbf{\texttt{\small DERIV:\allowbreak{}W6\_\allowbreak{}SOURCE\_\allowbreak{}GENERATOR\_\allowbreak{}SCORE\_\allowbreak{}FRAME:\allowbreak{}R10\_\allowbreak{}ORDER\_\allowbreak{}ZERO\_\allowbreak{}REDUCTION}} \hfill {\footnotesize conditional}\par
{\small Under the conditional-energy hypotheses H0-H4 (uniform q\_\allowbreak{}g-boundedness of P\_\allowbreak{}g, sup\_\allowbreak{}w K\textasciicircum{}0\_\allowbreak{}g <= kappa\_\allowbreak{}0 g\textasciicircum{}-2, J\_\allowbreak{}g <= kappa\_\allowbreak{}1 g\textasciicircum{}-2 (1+g\textasciicircum{}-2 E[V|w]), the fast-to-source conditional energy bound, and b\_\allowbreak{}g[m\_\allowbreak{}g] <= kappa\_\allowbreak{}3 g\textasciicircum{}-2), the complete R9 generator satisfies ||A(g)||\_\allowbreak{}(q\_\allowbreak{}g->q\_\allowbreak{}g) <= d\_\allowbreak{}0 g\textasciicircum{}-1 with the\,\ldots}\par
{\footnotesize \srcpath{docs/derivations/w6-source-generator-score-frame.md} --- \#\# 5. Exact reduction of the order-zero R10 bound (SF8-SF9)}\par
{\footnotesize \textit{Hypotheses.} H0-H4 with constants independent of g on 0<g<g\_\allowbreak{}*; each is stated explicitly in section 5.; The rank-one q\_\allowbreak{}g operator bound and R4 from the transport-budget note; 0<=e\_\allowbreak{}g<=E and gap gamma>0.}\par
{\footnotesize \textit{Remaining.} Prove H0-H4. The r=1,2 cases of R10 are not reduced;}\par
\smallskip
\noindent\textbf{\texttt{\small DERIV:\allowbreak{}WILSON\_\allowbreak{}BACKGROUND\_\allowbreak{}COVARIANCE\_\allowbreak{}CONTINUATION:\allowbreak{}BF1\_\allowbreak{}BF3}} \hfill {\footnotesize conditional}\par
{\small A common fast form with relative defect norm theta\_\allowbreak{}z<1 gives the energy-weighted inverse difference <=theta\_\allowbreak{}z/(1-theta\_\allowbreak{}z), the selected signed pairing bound BF1, and the exact energy-dependent Schur/source identity BF3.}\par
{\footnotesize \srcpath{docs/derivations/wilson-background-covariance-continuation.md} --- \#\# BF1. The energy-dependent relative-form statement (lines 48-102)}\par
{\footnotesize \textit{Hypotheses.} Finite periodic cubic SU(2) lattice with nondegenerate plaquettes and side L >= 3.; Product unit-S3 electric metric; epsilon > 0, k >= 0, lambda = k/epsilon; actual positive normalized Wilson ground Omega.; A0(z)=F0-z>=f\_\allowbreak{}z I>0; V is a bounded A0-relative form with norm <=theta\_\allowbreak{}z<1.}\par
{\footnotesize \textit{Remaining.} The normalized bounded inverse, selected pairing and Schur minimization are formalized. Identify the actual common closed form, inverse-square-root coordinates and compatible cross-functionals for the Wilson operator.}\par
\smallskip
\noindent\textbf{\texttt{\small DERIV:\allowbreak{}WILSON\_\allowbreak{}BACKGROUND\_\allowbreak{}COVARIANCE\_\allowbreak{}CONTINUATION:\allowbreak{}SC17}} \hfill {\footnotesize conditional}\par
{\small The exact distant mixed derivative is -R\_\allowbreak{}xy(F\_\allowbreak{}e tensor V\_\allowbreak{}f+V\_\allowbreak{}e tensor F\_\allowbreak{}f)-F\_\allowbreak{}e tensor F\_\allowbreak{}f; the original all-coupling proof does not supply a summable distance envelope.}\par
{\footnotesize \srcpath{docs/derivations/wilson-background-covariance-continuation.md} --- \#\# 5. The neutral channel and the exact remaining spatial bound (lines 494-534)}\par
{\footnotesize \textit{Hypotheses.} Finite periodic cubic SU(2) lattice with nondegenerate plaquettes and side L >= 3.; Product unit-S3 electric metric; epsilon > 0, k >= 0, lambda = k/epsilon; actual positive normalized Wilson ground Omega.; Distinct distant tails; V\_\allowbreak{}ef=0; R\_\allowbreak{}xy is the actual charged inverse.}\par
{\footnotesize \textit{Remaining.} Formalize this identity and the singlet obstruction; the all-coupling spatially summable estimate still needs its complete quantum reference-defect bound.}\par
\smallskip
\noindent\textbf{\texttt{\small DERIV:\allowbreak{}WILSON\_\allowbreak{}MARKED\_\allowbreak{}SHELL\_\allowbreak{}TRANSPORT:\allowbreak{}S18\_\allowbreak{}S20}} \hfill {\footnotesize conditional}\par
{\small The common weighted Taylor majorant sums finite-order Wilson/Hamiltonian matching; the relative shell gap stays >(t3/4)u\textasciicircum{}2 q(k) away from Gamma and carrier projections converge by S20.}\par
{\footnotesize \srcpath{docs/derivations/wilson-marked-shell-transport.md} --- \#\# 6. Sum the matching and transport the carrier (lines 391-461)}\par
{\footnotesize \textit{Hypotheses.} SU(3), fixed spatial spacing and small magnetic coupling on the S1 common domain.; Calibrated kinetic-window and the actual September 5 Hamiltonian G18 creator/activity/complete-band hypotheses and literal source charts are explicit inputs.; The Hamiltonian G18 gap and kinetic-window coefficient comparison apply on the same weighted literal-source frame; epsilon is sufficiently small.}\par
{\footnotesize \textit{Remaining.} Formalize power-series tail convergence, the common polar source chart, relative spectral perturbation and exact symmetry cancellation at Gamma.}\par
\smallskip
\noindent\textbf{\texttt{\small DERIV:\allowbreak{}WILSON\_\allowbreak{}SC17\_\allowbreak{}SPATIAL\_\allowbreak{}CLOSURE:\allowbreak{}PROPOSED\_\allowbreak{}QUANTUM\_\allowbreak{}PARAMETER\_\allowbreak{}WINDOW}} \hfill {\footnotesize conditional}\par
{\small Under the proposed numerical reference and normalized-defect inputs, the scalar Riccati parameter satisfies p(lambda)<=p(3/100)=96228*sqrt(3)/625000<27/100<1, and the small root satisfies 0<=r\_\allowbreak{}minus<a\_\allowbreak{}min throughout 0<=lambda<=3/100. These are conditional parameter conclusions; the section's complete Wilson defect estimate and exponential spatial decay are not established by this arithmetic.}\par
{\footnotesize \srcpath{docs/derivations/wilson-sc17-spatial-closure.md} --- \#\# 6. Formal repair of the spatial closure for lambda >= 1/73 via the quantum reference defect (lines 781-821)}\par
{\footnotesize \textit{Hypotheses.} Real coupling 0<=lambda<=3/100; epsilon>0 when interpreting the normalized defect c=delta/(2epsilon).; The companion script's proposed defaults are L=3, C\_\allowbreak{}s=6/5 and reference velocity one; beta=(108/25)*sqrt(33), c=lambda*sqrt(lambda)/48 and a\_\allowbreak{}min=1/(3*sqrt(33)).; p(lambda)=4*beta\textasciicircum{}2*c and r\_\allowbreak{}minus=(1-sqrt(1-p(lambda)))/(2*beta).; These numerical inputs are supplied parameters. Their identification with a uniform semigroup budget, the complete conditional quantum-pressure defect and the actual reference floor is a separate model obligation, not a consequence of the quartic Taylor coefficient.}\par
{\footnotesize \textit{Remaining.} The scalar comparison is proved in the supported successor SC17\_\allowbreak{}P. To apply it to the source section's spatial-closure claim, establish the actual weighted reference budget and complete outside-pressure, metric, projector and cutoff defect bound, derive the ground-state Riccati comparison, and justify an exponential spatial weight.}\par
\smallskip
\noindent\textbf{\texttt{\small DERIV:\allowbreak{}WILSON\_\allowbreak{}SC17\_\allowbreak{}SPATIAL\_\allowbreak{}CLOSURE:\allowbreak{}R10\_\allowbreak{}R11}} \hfill {\footnotesize conditional}\par
{\small Uniform block Hessian floors 2a\_\allowbreak{}B,2a\_\allowbreak{}C and cross norm 2h\_\allowbreak{}BC yield c\_\allowbreak{}BC<=h\_\allowbreak{}BC/sqrt(a\_\allowbreak{}B a\_\allowbreak{}C); a reference angle budget plus full Hessian error r yields physical gap >=2epsilon[a0(1-kappa\_\allowbreak{}A)-2r].}\par
{\footnotesize \srcpath{docs/derivations/wilson-sc17-spatial-closure.md} --- \#\# 3. An explicit true-angle consequence if the full Hessian defect is known (lines 649-697)}\par
{\footnotesize \textit{Hypotheses.} Finite periodic cubic SU(2) lattice with nondegenerate plaquettes and side L >= 3.; Product unit-S3 electric metric; epsilon > 0, k >= 0, lambda = k/epsilon; actual positive normalized Wilson ground Omega.; a\_\allowbreak{}B,a\_\allowbreak{}C>0; h\_\allowbreak{}BC\textasciicircum{}2<a\_\allowbreak{}B a\_\allowbreak{}C; kappa\_\allowbreak{}A<1; 2r<a0(1-kappa\_\allowbreak{}A).}\par
{\footnotesize \textit{Remaining.} Formalize Brascamp-Lieb conditional variance, marginal Schur curvature and the actual projection-angle assembly.}\par
\smallskip
\noindent\textbf{\texttt{\small DERIV:\allowbreak{}WILSON\_\allowbreak{}SC17\_\allowbreak{}SPATIAL\_\allowbreak{}CLOSURE:\allowbreak{}R12\_\allowbreak{}R14}} \hfill {\footnotesize conditional}\par
{\small The necessary full conditional Hessian defect is Hess(V\_\allowbreak{}B+((H\_\allowbreak{}out Omega)/Omega))-2epsilon A\textasciicircum{}2; exact Gaussian pressure cancels the reference BB* term, while cutoffs create the explicit gradient and commutator defects.}\par
{\footnotesize \srcpath{docs/derivations/wilson-sc17-spatial-closure.md} --- \#\# 4. The actual quantum pressure and cutoff defects are part of the input (lines 698-756)}\par
{\footnotesize \textit{Hypotheses.} Finite periodic cubic SU(2) lattice with nondegenerate plaquettes and side L >= 3.; Product unit-S3 electric metric; epsilon > 0, k >= 0, lambda = k/epsilon; actual positive normalized Wilson ground Omega.; Actual vacuum and compatible Euclidean block chart; all multiplication terms assigned once.}\par
{\footnotesize \textit{Remaining.} Formalize the outside-pressure derivative, Gaussian cancellation and every cutoff/product-domain term; the actual uniform smallness estimate remains unproved.}\par
\smallskip
\noindent\textbf{\texttt{\small DERIV:\allowbreak{}WILSON\_\allowbreak{}SC17\_\allowbreak{}SPATIAL\_\allowbreak{}CLOSURE:\allowbreak{}R4\_\allowbreak{}R9}} \hfill {\footnotesize conditional}\par
{\small If beta\_\allowbreak{}s(A)<infinity and the full conditional potential Hessian defect is <=delta with 2 beta\_\allowbreak{}s(A)\textasciicircum{}2 delta/epsilon<1, the true log-ground Hessian error is <=r\_\allowbreak{}minus and conditional gap >=2epsilon(a-r\_\allowbreak{}minus).}\par
{\footnotesize \srcpath{docs/derivations/wilson-sc17-spatial-closure.md} --- \#\# 2. A rigorous conditional small-defect criterion (lines 574-648)}\par
{\footnotesize \textit{Hypotheses.} Finite periodic cubic SU(2) lattice with nondegenerate plaquettes and side L >= 3.; Product unit-S3 electric metric; epsilon > 0, k >= 0, lambda = k/epsilon; actual positive normalized Wilson ground Omega.; A>=aI>0; the full conditional potential includes outside quantum pressure; smooth ground convergence and domain conditions in the section hold.}\par
{\footnotesize \textit{Remaining.} Formalize the noncommutative weighted norm, Duhamel Riccati comparison, ground limit and curvature-to-gap passage.}\par
\smallskip
\noindent\textbf{\texttt{\small DERIV:\allowbreak{}WILSON\_\allowbreak{}SELECTED\_\allowbreak{}INVERSE\_\allowbreak{}WALL:\allowbreak{}W4\_\allowbreak{}W5}} \hfill {\footnotesize conditional}\par
{\small On admissible reference graph vectors, S\_\allowbreak{}g-S\_\allowbreak{}0=A\_\allowbreak{}g-r\_\allowbreak{}g*R\_\allowbreak{}g r\_\allowbreak{}g; only the selected inverse bound W5 plus direct/force remainders is sufficient for the displayed O(g\textasciicircum{}3) comparison.}\par
{\footnotesize \srcpath{docs/derivations/wilson-selected-inverse-wall.md} --- \#\# 3. Repair: only compare inverses on the graph force (lines 101-157)}\par
{\footnotesize \textit{Hypotheses.} The operator chart, vacuum subtraction and fast/retained spaces are those of the spatial Schur derivation; domains are explicitly checked rather than identified by notation.; z lies below both full fast spectra; J\_\allowbreak{}z p and the interacting minimizer are admissible; W5 includes f>0,a3,t2,t0,c\_\allowbreak{}sel.}\par
{\footnotesize \textit{Remaining.} The bounded Hilbert-space completion/minimum and weighted pairing are formalized. Construct the actual unbounded fast-form domain, source identification and reduced closed-form Schur operator.}\par
\smallskip
\noindent\textbf{\texttt{\small DERIV:\allowbreak{}WILSON\_\allowbreak{}SPATIAL\_\allowbreak{}SCHUR\_\allowbreak{}EXCESS:\allowbreak{}SP1\_\allowbreak{}SP6}} \hfill {\footnotesize conditional}\par
{\small The exact full-reference Schur excess is A\_\allowbreak{}g-V\_\allowbreak{}g*(I+C\_\allowbreak{}g)\textasciicircum{}-1 V\_\allowbreak{}g; the minimizing graph, two-sided form bounds and -dS\_\allowbreak{}g/dz=J\_\allowbreak{}g,z*J\_\allowbreak{}g,z retain the induced metric.}\par
{\footnotesize \srcpath{docs/derivations/wilson-spatial-schur-excess.md} --- \#\# 2. Exact excess about the full reference graph (lines 42-114)}\par
{\footnotesize \textit{Hypotheses.} Closed self-adjoint reference and interacting forms on the stated compatible P/Q domains, already expressed in a common physical source chart and vacuum-subtracted.; Real energy z below the full fast restriction; positive retained energy weight B.; F\_\allowbreak{}z>=f\_\allowbreak{}zI>0; ||C\_\allowbreak{}g||<=theta<1; all weighted A\_\allowbreak{}g,V\_\allowbreak{}g extend boundedly.}\par
{\footnotesize \textit{Remaining.} The normalized inverse, positive fast form and unique bounded Schur minimum are formalized. Identify actual Wilson form domains and sources;}\par
\smallskip
\noindent\textbf{\texttt{\small DERIV:\allowbreak{}WILSON\_\allowbreak{}SPATIAL\_\allowbreak{}SCHUR\_\allowbreak{}EXCESS:\allowbreak{}SP20\_\allowbreak{}SP24}} \hfill {\footnotesize conditional}\par
{\small Complete Schur comparisons give the reciprocal-gap recursion, product/inverse-fast-energy lower bound SP21-SP22 and source-overlap telescoping SP24 in common physical clock units.}\par
{\footnotesize \srcpath{docs/derivations/wilson-spatial-schur-excess.md} --- \#\# 6. Iterating energies and observable amplitudes in the physical clock (lines 412-477)}\par
{\footnotesize \textit{Hypotheses.} Closed self-adjoint reference and interacting forms on the stated compatible P/Q domains, already expressed in a common physical source chart and vacuum-subtracted.; Real energy z below the full fast restriction; positive retained energy weight B.; 0<alpha\_\allowbreak{}j<=1; full fast floors f\_\allowbreak{}j; summable epsilon\_\allowbreak{}j=1-alpha\_\allowbreak{}j bounded below 1; summable inverse fast floors; actual source maps and errors eta\_\allowbreak{}j.}\par
{\footnotesize \textit{Remaining.} Formalize iteration of the closed-form gap theorem, positive infinite products and complete carrier/source projection transport.}\par
\smallskip
\noindent\textbf{\texttt{\small DERIV:\allowbreak{}WILSON\_\allowbreak{}SPATIAL\_\allowbreak{}SCHUR\_\allowbreak{}EXCESS:\allowbreak{}SP7\_\allowbreak{}SP10}} \hfill {\footnotesize conditional}\par
{\small K1=J\_\allowbreak{}z*W1J\_\allowbreak{}z and K2=J\_\allowbreak{}z*W2J\_\allowbreak{}z-V1*V1 are the complete Schur coefficients; SP8 gives the explicit O(|g|\textasciicircum{}3) remainder SP9, including the dressed force correction SP10.}\par
{\footnotesize \srcpath{docs/derivations/wilson-spatial-schur-excess.md} --- \#\# 3. The complete second-order coefficient and a uniform remainder (lines 115-182)}\par
{\footnotesize \textit{Hypotheses.} Closed self-adjoint reference and interacting forms on the stated compatible P/Q domains, already expressed in a common physical source chart and vacuum-subtracted.; Real energy z below the full fast restriction; positive retained energy weight B.; SP8 remainder constants a3,v2,v1,c1,g0 are finite and c1|g|<=theta<1.}\par
{\footnotesize \textit{Remaining.} The genuine Neumann operator expansion and cubic norm/sandwich remainder are formalized. Assemble the g-dependent direct term, moving forces, fast defect and spatial coefficients on the actual compatible Wilson form domain.}\par
\smallskip
\noindent\textbf{\texttt{\small DERIV:\allowbreak{}WILSON\_\allowbreak{}TRUE\_\allowbreak{}VACUUM\_\allowbreak{}BLOCK\_\allowbreak{}ESTIMATES:\allowbreak{}BA9\_\allowbreak{}BA10}} \hfill {\footnotesize conditional}\par
{\small The true conditional operator includes outside pressure W\_\allowbreak{}B=(H\_\allowbreak{}out Omega)/Omega and has gap >=gap(h\_\allowbreak{}B(b))-osc\_\allowbreak{}B W\_\allowbreak{}B; a frozen-block gap alone omits this term.}\par
{\footnotesize \srcpath{docs/derivations/wilson-true-vacuum-block-estimates.md} --- \#\# BA4. What a frozen-block spectral computation would additionally require (lines 142-168)}\par
{\footnotesize \textit{Hypotheses.} Finite periodic cubic SU(2) lattice with nondegenerate plaquettes and side L >= 3.; Product unit-S3 electric metric; epsilon > 0, k >= 0, lambda = k/epsilon; actual positive normalized Wilson ground Omega.; All conditional blocks use the actual vacuum law mu=Omega\textasciicircum{}2dU, not isolated-block spectral data.; The frozen-block gap and finite pressure oscillation are available on compatible form domains.}\par
{\footnotesize \textit{Remaining.} Formalize the conditional ground transform and bounded-potential min-max comparison, including the domain of H\_\allowbreak{}out Omega.}\par
\smallskip
\noindent\textbf{\texttt{\small DERIV:\allowbreak{}WILSON\_\allowbreak{}VACUUM\_\allowbreak{}ALIGNED\_\allowbreak{}ASSEMBLY:\allowbreak{}VA12\_\allowbreak{}VA13}} \hfill {\footnotesize conditional}\par
{\small For physical conditional-angle row kappa<1, G\_\allowbreak{}mu\textasciicircum{}2>=(1-kappa)G\_\allowbreak{}mu and gap\_\allowbreak{}phys(H)>=gamma\_\allowbreak{}*(1-kappa); the single plaquette physical restriction has E\_\allowbreak{}e=E0 and zero pair angle.}\par
{\footnotesize \srcpath{docs/derivations/wilson-vacuum-aligned-assembly.md} --- \#\# VA5. Projection geometry, with the Gauss law retained (lines 211-251)}\par
{\footnotesize \textit{Hypotheses.} Finite periodic cubic SU(2) lattice with nondegenerate plaquettes and side L >= 3.; Product unit-S3 electric metric; epsilon > 0, k >= 0, lambda = k/epsilon; actual positive normalized Wilson ground Omega.; c\_\allowbreak{}ef\textasciicircum{}phys is the actual L2 conditional-projection norm, and sup\_\allowbreak{}e sum\_\allowbreak{}(f!=e)c\_\allowbreak{}ef\textasciicircum{}phys<=kappa<1.}\par
{\footnotesize \textit{Remaining.} The arbitrary-dimensional bounded projection-sum implication and kernel-complement gap are formalized. Construct actual conditional projections, derive the local anticommutator bound from their Friedrichs angles, identify the common kernel as constants and establish the required row budget below one.}\par
\smallskip
\noindent\textbf{\texttt{\small DERIV:\allowbreak{}WILSON\_\allowbreak{}VACUUM\_\allowbreak{}ALIGNED\_\allowbreak{}ASSEMBLY:\allowbreak{}VA14\_\allowbreak{}VA19}} \hfill {\footnotesize conditional}\par
{\small Approximate tensorization or a summable block-angle row gives the stated physical gap with covering multiplicity; the original VA10 local floor tends to zero on the displayed logarithmic continuum trajectory.}\par
{\footnotesize \srcpath{docs/derivations/wilson-vacuum-aligned-assembly.md} --- \#\# VA6. The next estimate and the attempts to establish it (lines 252-415)}\par
{\footnotesize \textit{Hypotheses.} Finite periodic cubic SU(2) lattice with nondegenerate plaquettes and side L >= 3.; Product unit-S3 electric metric; epsilon > 0, k >= 0, lambda = k/epsilon; actual positive normalized Wilson ground Omega.; Actual conditional blocks, gamma\_\allowbreak{}min>0 and finite covering multiplicity m\_\allowbreak{}*; a spatial envelope is an additional hypothesis, not assigned data.}\par
{\footnotesize \textit{Remaining.} Actual-measure covariance minorization and explicit-domain unitary gap transport are formalized. Establish the actual Wilson conditional density floor, conditional disintegration, covariance-to-projection-angle norm and spatial row-budget bounds.}\par
\smallskip
\noindent\textbf{\texttt{\small DERIV:\allowbreak{}WILSON\_\allowbreak{}WEIGHTED\_\allowbreak{}REPAIR\_\allowbreak{}AND\_\allowbreak{}ROTOR\_\allowbreak{}GAP:\allowbreak{}WR17\_\allowbreak{}WR21}} \hfill {\footnotesize conditional}\par
{\small The actual trial comparison needs gamma\_\allowbreak{}Lambda-(mean R-inf R)>0 uniformly; a gauge-invariant patch has full-boundary trial excitations of energy 3epsilon, so an unrestricted boundary fast floor of order sqrt(epsilon k) cannot be inferred.}\par
{\footnotesize \srcpath{docs/derivations/wilson-weighted-repair-and-rotor-gap.md} --- \#\# WR6. The precise remaining derivation (lines 229-308)}\par
{\footnotesize \textit{Hypotheses.} Each operator is identified separately: classical quotient diffusion, complete compact S3 quantum rotor, actual coupled link lattice or independent product model as stated in the section.; Positive normalized overlapping-plaquette trial; gamma\_\allowbreak{}Lambda is an actual weighted generator gap.}\par
{\footnotesize \textit{Remaining.} Formalize the common form domains and mean-residual comparison; construct Haar-link marginal trial vectors proving the full-boundary obstruction.}\par
\smallskip
\noindent\textbf{\texttt{\small DERIV:\allowbreak{}YANGMILLS\_\allowbreak{}CONTINUUM\_\allowbreak{}BALABAN\_\allowbreak{}MULTISCALE\_\allowbreak{}PROOF:\allowbreak{}GHOST\_\allowbreak{}5\_\allowbreak{}8}} \hfill {\footnotesize conditional}\par
{\small Rescaling z=g*zeta exposes g\textasciicircum{}(1-kappa\_\allowbreak{}prime) in (5.8), but the required volume/scale-uniform covariant inverse-gradient bound is not proved by that algebra.}\par
{\footnotesize \srcpath{docs/derivations/yangmills-continuum-balaban-multiscale-proof.md} --- \#\#\# 5.2 Faddeev--Popov Ghost Determinant (lines 198-208)}\par
{\footnotesize \textit{Hypotheses.} Submitted SU(N), N>=2 continuum multiscale construction; read together with its in-document Repair record audited 2026-09-09.; The source title and summary claim completion, but its repair record retains explicit averaging/interpolation, irrelevance, physical-scale and scattering obligations. No machine certification is inferred from prose.; 0<kappa\_\allowbreak{}prime<1/2; a||zeta||\_\allowbreak{}infinity<=g\textasciicircum{}-kappa\_\allowbreak{}prime; the first operator inequality in (5.8) is an additional input.}\par
{\footnotesize \textit{Remaining.} Construct the precise ghost operator, determinant class and uniform inverse-gradient estimate on its constrained domain; scalar rescaling alone does not discharge them.}\par
\smallskip
\noindent\textbf{\texttt{\small DERIV:\allowbreak{}YANGMILLS\_\allowbreak{}CONTINUUM\_\allowbreak{}BALABAN\_\allowbreak{}MULTISCALE\_\allowbreak{}PROOF:\allowbreak{}OS0}} \hfill {\footnotesize conditional}\par
{\small The asserted Schwinger growth bound (8.2) supplies OS regularity only after its actual uniform moment/Sobolev estimates and continuum fields are constructed.}\par
{\footnotesize \srcpath{docs/derivations/yangmills-continuum-balaban-multiscale-proof.md} --- \#\#\# 8.1 OS0: Analyticity and Temperedness (lines 368-374)}\par
{\footnotesize \textit{Hypotheses.} Submitted SU(N), N>=2 continuum multiscale construction; read together with its in-document Repair record audited 2026-09-09.; The source title and summary claim completion, but its repair record retains explicit averaging/interpolation, irrelevance, physical-scale and scattering obligations. No machine certification is inferred from prose.}\par
{\footnotesize \textit{Remaining.} Formalize the continuum Schwinger distributions and prove the factorial/test-seminorm bound from the actual regulated measures.}\par
\smallskip
\noindent\textbf{\texttt{\small DERIV:\allowbreak{}YANGMILLS\_\allowbreak{}CONTINUUM\_\allowbreak{}BALABAN\_\allowbreak{}MULTISCALE\_\allowbreak{}PROOF:\allowbreak{}OS3}} \hfill {\footnotesize conditional}\par
{\small Bosonic multiplication observables have permutation-symmetric Schwinger functions, provided the claimed continuum distributions exist.}\par
{\footnotesize \srcpath{docs/derivations/yangmills-continuum-balaban-multiscale-proof.md} --- \#\#\# 8.4 OS3: Permutation Symmetry (lines 409-411)}\par
{\footnotesize \textit{Hypotheses.} Submitted SU(N), N>=2 continuum multiscale construction; read together with its in-document Repair record audited 2026-09-09.; The source title and summary claim completion, but its repair record retains explicit averaging/interpolation, irrelevance, physical-scale and scattering obligations. No machine certification is inferred from prose.}\par
{\footnotesize \textit{Remaining.} Formalize products of the constructed fields/observables and the permutation action on their distributional arguments.}\par
\smallskip
\noindent\textbf{\texttt{\small DERIV:\allowbreak{}YANGMILLS\_\allowbreak{}CONTINUUM\_\allowbreak{}BALABAN\_\allowbreak{}MULTISCALE\_\allowbreak{}PROOF:\allowbreak{}OS\_\allowbreak{}RECONSTRUCTION}} \hfill {\footnotesize conditional}\par
{\small With corrected OS hypotheses and totality, the reflection-positive quotient completion has a physical Hilbert space, vacuum and nonnegative self-adjoint time generator.}\par
{\footnotesize \srcpath{docs/derivations/yangmills-continuum-balaban-multiscale-proof.md} --- \#\#\# 9.1 The Physical Hilbert Space and Hamiltonian (lines 439-453)}\par
{\footnotesize \textit{Hypotheses.} Submitted SU(N), N>=2 continuum multiscale construction; read together with its in-document Repair record audited 2026-09-09.; The source title and summary claim completion, but its repair record retains explicit averaging/interpolation, irrelevance, physical-scale and scattering obligations. No machine certification is inferred from prose.}\par
{\footnotesize \textit{Remaining.} Formalize the OS quotient, completion, translations and corrected reconstruction regularity; actual OS inputs remain separate obligations.}\par
\smallskip
\noindent\textbf{\texttt{\small DERIV:\allowbreak{}YANGMILLS\_\allowbreak{}CONTINUUM\_\allowbreak{}BALABAN\_\allowbreak{}MULTISCALE\_\allowbreak{}PROOF:\allowbreak{}SCATTERING}} \hfill {\footnotesize conditional}\par
{\small Haag-Ruelle scattering is conditional on an isolated positive single-particle mass shell; the source now explicitly states that dispersion hypothesis. Nonzero off-shell/tree-level four-point data alone do not supply the required on-shell amplitude lower bound.}\par
{\footnotesize \srcpath{docs/derivations/yangmills-continuum-balaban-multiscale-proof.md} --- \#\#\# 9.3 Relativistic Wightman Fields and Non-Trivial Scattering (\$S \textbackslash{}ne I\$) (lines 479-501)}\par
{\footnotesize \textit{Hypotheses.} Submitted SU(N), N>=2 continuum multiscale construction; read together with its in-document Repair record audited 2026-09-09.; The source title and summary claim completion, but its repair record retains explicit averaging/interpolation, irrelevance, physical-scale and scattering obligations. No machine certification is inferred from prose.; Isolated mass hyperboloid separated below multiparticle threshold; locality, spectrum and actual LSZ/Haag-Ruelle assumptions.}\par
{\footnotesize \textit{Remaining.} Formalize the isolated-shell hypothesis, asymptotic states and a controlled on-shell nonzero amplitude; establish these for the constructed theory before concluding S!=I.}\par
\smallskip
\noindent\textbf{\texttt{\small DERIV:\allowbreak{}YANGMILLS\_\allowbreak{}CONTINUUM\_\allowbreak{}BALABAN\_\allowbreak{}MULTISCALE\_\allowbreak{}PROOF:\allowbreak{}THEOREM\_\allowbreak{}4\_\allowbreak{}2}} \hfill {\footnotesize conditional}\par
{\small The large-field damping claim (4.4)-(4.5) is conditional on a positive coarse-action subtraction constant and the unproved irrelevance gain identified in the repair record.}\par
{\footnotesize \srcpath{docs/derivations/yangmills-continuum-balaban-multiscale-proof.md} --- \#\#\# 4.2 Super-Exponential Suppression of Large Fields (lines 157-180)}\par
{\footnotesize \textit{Hypotheses.} Submitted SU(N), N>=2 continuum multiscale construction; read together with its in-document Repair record audited 2026-09-09.; The source title and summary claim completion, but its repair record retains explicit averaging/interpolation, irrelevance, physical-scale and scattering obligations. No machine certification is inferred from prose.; 0<kappa\_\allowbreak{}prime<kappa<1/6; actual regularized functional measure and corrected background construction; uniform constants.}\par
{\footnotesize \textit{Remaining.} Prove the quantitative irrelevance/subtraction constant and the actual large-field measure estimate, including the averaging constraint and fluctuation direction invisible to a coarse penalty.}\par
\smallskip
\noindent\textbf{\texttt{\small DERIV:\allowbreak{}YANGMILLS\_\allowbreak{}CONTINUUM\_\allowbreak{}BALABAN\_\allowbreak{}MULTISCALE\_\allowbreak{}PROOF:\allowbreak{}THEOREM\_\allowbreak{}6\_\allowbreak{}2}} \hfill {\footnotesize conditional}\par
{\small The claimed RG contraction (6.5) has the repaired induction envelope ||K\_\allowbreak{}k||<=4 C\_\allowbreak{}ind g\_\allowbreak{}k\textasciicircum{}2, but its single-block contraction gain still requires the irrelevance lemma.}\par
{\footnotesize \srcpath{docs/derivations/yangmills-continuum-balaban-multiscale-proof.md} --- \#\#\# 6.3 The Mayer Tree-Graph Expansion \& Cluster Contraction (lines 269-295)}\par
{\footnotesize \textit{Hypotheses.} Submitted SU(N), N>=2 continuum multiscale construction; read together with its in-document Repair record audited 2026-09-09.; The source title and summary claim completion, but its repair record retains explicit averaging/interpolation, irrelevance, physical-scale and scattering obligations. No machine certification is inferred from prose.; lambda\_\allowbreak{}contraction in (0,1/2); sufficiently small initial g0; the corrected averaging, fluctuation and large-field inputs hold.}\par
{\footnotesize \textit{Remaining.} Formalize the actual tree expansion and prove the residual local-operator irrelevance gain after marginal subtraction; then prove the model-dependent contraction.}\par
\smallskip
\noindent\textbf{\texttt{\small DERIV:\allowbreak{}YANGMILLS\_\allowbreak{}CONTINUUM\_\allowbreak{}BALABAN\_\allowbreak{}MULTISCALE\_\allowbreak{}PROOF:\allowbreak{}THEOREM\_\allowbreak{}6\_\allowbreak{}3}} \hfill {\footnotesize conditional}\par
{\small The source records b0=11N/(48pi\textasciicircum{}2), the perturbative running formula (6.11)-(6.12) and a uniform polymer conclusion (Corollary 6.4); its repaired induction requires quantitative remainder control and a propagated envelope.}\par
{\footnotesize \srcpath{docs/derivations/yangmills-continuum-balaban-multiscale-proof.md} --- \#\#\# 6.4 Asymptotic Freedom and Infinite Scale Stability (lines 296-314)}\par
{\footnotesize \textit{Hypotheses.} Submitted SU(N), N>=2 continuum multiscale construction; read together with its in-document Repair record audited 2026-09-09.; The source title and summary claim completion, but its repair record retains explicit averaging/interpolation, irrelevance, physical-scale and scattering obligations. No machine certification is inferred from prose.; Perturbative UV regime; controlled beta remainder and actual RG map; this does not place the hadronic scale in the small-g domain.}\par
{\footnotesize \textit{Remaining.} Formalize the actual RG beta coefficient and remainder assumptions before asserting monotonicity for all g0; prove the stated envelope with the repaired constants.}\par
\smallskip
\noindent\textbf{\texttt{\small DERIV:\allowbreak{}YANGMILLS\_\allowbreak{}CONTINUUM\_\allowbreak{}BALABAN\_\allowbreak{}MULTISCALE\_\allowbreak{}PROOF:\allowbreak{}THEOREM\_\allowbreak{}7\_\allowbreak{}2}} \hfill {\footnotesize conditional}\par
{\small The continuum measure claim requires a continuous positive-definite characteristic functional on the full nuclear test space; convergence only of gauge-invariant observables is not itself this construction.}\par
{\footnotesize \srcpath{docs/derivations/yangmills-continuum-balaban-multiscale-proof.md} --- \#\#\# 7.2 Construction of the Limiting Continuum Measure \$d\textbackslash{}mu\_\allowbreak{}\{\textbackslash{}text\{YM\}\}\$ (lines 341-363)}\par
{\footnotesize \textit{Hypotheses.} Submitted SU(N), N>=2 continuum multiscale construction; read together with its in-document Repair record audited 2026-09-09.; The source title and summary claim completion, but its repair record retains explicit averaging/interpolation, irrelevance, physical-scale and scattering obligations. No machine certification is inferred from prose.; Theorem 7.1 with the repaired summable bound; actual uniform regularity and characteristic-functional estimates.}\par
{\footnotesize \textit{Remaining.} Construct the characteristic functional with continuity/positivity and a consistent gauge-invariant observable realization, then apply Minlos-Bochner.}\par
\smallskip
\noindent\textbf{\texttt{\small DERIV:\allowbreak{}YANGMILLS\_\allowbreak{}CONTINUUM\_\allowbreak{}BALABAN\_\allowbreak{}MULTISCALE\_\allowbreak{}PROOF:\allowbreak{}THEOREM\_\allowbreak{}8\_\allowbreak{}1}} \hfill {\footnotesize conditional}\par
{\small Centered reflection-adapted deterministic blocking preserves reflection positivity under its exact support/equivariance hypotheses, and a convergent compatible limit retains the inequality (8.5).}\par
{\footnotesize \srcpath{docs/derivations/yangmills-continuum-balaban-multiscale-proof.md} --- \#\#\# 8.3 OS2: Reflection Positivity (lines 383-408)}\par
{\footnotesize \textit{Hypotheses.} Submitted SU(N), N>=2 continuum multiscale construction; read together with its in-document Repair record audited 2026-09-09.; The source title and summary claim completion, but its repair record retains explicit averaging/interpolation, irrelevance, physical-scale and scattering obligations. No machine certification is inferred from prose.; Reflection-adapted deterministic block map, positive-time support and actual convergent correlations; the cited centered-blocking construction applies.}\par
{\footnotesize \textit{Remaining.} Formalize the deterministic reflected blocking and limit passage for the actual continuum measure; general reflection-equivariant Markov blocking is not this theorem.}\par
\smallskip
\noindent\textbf{\texttt{\small DERIV:\allowbreak{}YANGMILLS\_\allowbreak{}CONTINUUM\_\allowbreak{}BALABAN\_\allowbreak{}MULTISCALE\_\allowbreak{}PROOF:\allowbreak{}THEOREM\_\allowbreak{}9\_\allowbreak{}1}} \hfill {\footnotesize conditional}\par
{\small A common strictly positive exponential physical-time correlation rate on the dense centered reconstructed family implies spectrum(H) subset \{0\} union [Delta\_\allowbreak{}phys,infinity), with the simple vacuum; this conditional spectral argument is confirmed sound by the source repair record.}\par
{\footnotesize \srcpath{docs/derivations/yangmills-continuum-balaban-multiscale-proof.md} --- \#\#\# 9.2 The Physical Spectral Mass Gap (lines 454-478)}\par
{\footnotesize \textit{Hypotheses.} Submitted SU(N), N>=2 continuum multiscale construction; read together with its in-document Repair record audited 2026-09-09.; The source title and summary claim completion, but its repair record retains explicit averaging/interpolation, irrelevance, physical-scale and scattering obligations. No machine certification is inferred from prose.; Actual reconstructed physical space with dense observable family; common rate Delta\_\allowbreak{}phys>0 in physical time; simple vacuum.}\par
{\footnotesize \textit{Remaining.} Formalize the spectral-measure positivity, infinite-time exclusion and density extension; derive the actual physical clustering input (9.7) separately.}\par
\smallskip
\noindent\textbf{\texttt{\small DERIV:\allowbreak{}YANGMILLS\_\allowbreak{}GPU\_\allowbreak{}RESOLUTION\_\allowbreak{}AUDIT:\allowbreak{}GA8\_\allowbreak{}GA9}} \hfill {\footnotesize conditional}\par
{\small The inserted scalar recursion has explicit positive product/tail bounds GA8 when its assumptions hold; the Phase 4 four-plaquette builder remains a tensor sum GA9.}\par
{\footnotesize \srcpath{docs/derivations/yangmills-gpu-resolution-audit.md} --- \#\# GA4. The additional Phase 4 has a valid conditional scalar limit (lines 162-209)}\par
{\footnotesize \textit{Hypotheses.} The preserved September 9 scripts are inspected as the operators they actually implement; they are not assumed to represent Wilson dynamics.; Analytic matrix identities, replayed floating observations and reported-only runs remain distinct.; eps\_\allowbreak{}j=C(ell+jh)\textasciicircum{}(-3/2)<1; positive initial gap and fast floors; these are supplied inputs in the script.}\par
{\footnotesize \textit{Remaining.} Formalize the scalar positive-product and geometric tail bounds; derive actual coarse Wilson parameters before applying them.}\par
\smallskip
\noindent\textbf{\texttt{\small DERIV:\allowbreak{}YANGMILLS\_\allowbreak{}RECONSTRUCTION:\allowbreak{}R6\_\allowbreak{}R7}} \hfill {\footnotesize conditional}\par
{\small If C\_\allowbreak{}F(t)<=A\_\allowbreak{}F exp(-eta*t+alpha*R)+B\_\allowbreak{}F exp(-beta*R) for every R,t>=0, choose R=eta*t/(alpha+beta) to obtain common rate eta*beta/(alpha+beta).}\par
{\footnotesize \srcpath{docs/derivations/yangmills-reconstruction.md} --- \#\# 5. A precise sufficient repair, with its remaining hypothesis exposed (lines 90-111)}\par
{\footnotesize \textit{Hypotheses.} A nonnegative self-adjoint physical Hamiltonian with a specified vacuum and its positive spectral measures; physical Euclidean time is used.; OS reconstruction and any limit theorem are explicit analytic inputs, not consequences of the finite controls.; eta,beta>0,alpha>=0; one common rate and controlled cutoff/volume constants on a dense centered physical family.}\par
{\footnotesize \textit{Remaining.} Real-radius exponent balancing, actual-measure spectral exclusion and dense-family extension are formalized. Establish the actual uniform Yang-Mills two-parameter localization estimate, physical spectral identification and total family;}\par
\smallskip
\noindent\textbf{\texttt{\small DERIV:\allowbreak{}YANGMILLS\_\allowbreak{}RESOLVENT\_\allowbreak{}LOCALIZATION\_\allowbreak{}AND\_\allowbreak{}DIRICHLET\_\allowbreak{}GAP:\allowbreak{}LOCAL\_\allowbreak{}KERNEL\_\allowbreak{}IMPLICATION}} \hfill {\footnotesize conditional}\par
{\small Given uniform massive exponentially decaying kernels and a finite-range perturbation of size C2|g|, the localized scalar pairing has a volume-independent geometric-sum bound.}\par
{\footnotesize \srcpath{docs/derivations/yangmills-resolvent-localization-and-dirichlet-gap.md} --- \#\#\# 2.3 Volume-Uniformity of the Pairing (lines 62-80)}\par
{\footnotesize \textit{Hypotheses.} This submitted proposal is explicitly preserved under its September 9 source audit; closure claims are not the current graph verdict.; The actual model/operator identification is an obligation, separate from a scalar finite replay.; Kernel decay constants and perturbation range/row bounds uniform in volume; finite localized source with the source normalization explicitly retained.}\par
{\footnotesize \textit{Remaining.} Infinite kernel-pairing summability, exact integer-lattice sums and the full source l1 factor are formalized. Identify actual Wilson fast kernels and connected signed sources and prove the required uniform decay and perturbation envelopes.}\par
\smallskip
\noindent\textbf{\texttt{\small DERIV:\allowbreak{}YANGMILLS\_\allowbreak{}RESOLVENT\_\allowbreak{}LOCALIZATION\_\allowbreak{}AND\_\allowbreak{}DIRICHLET\_\allowbreak{}GAP:\allowbreak{}RESOLVENT\_\allowbreak{}IDENTITY}} \hfill {\footnotesize conditional}\par
{\small Compatible invertible fast operators obey R\_\allowbreak{}g-R0=-R0(F\_\allowbreak{}g-F0)R\_\allowbreak{}g and the stated selected-pairing identity.}\par
{\footnotesize \srcpath{docs/derivations/yangmills-resolvent-localization-and-dirichlet-gap.md} --- \#\#\# 2.1 The Resolvent Identity (lines 42-54)}\par
{\footnotesize \textit{Hypotheses.} This submitted proposal is explicitly preserved under its September 9 source audit; closure claims are not the current graph verdict.; The actual model/operator identification is an obligation, separate from a scalar finite replay.; Real z below both spectra; compatible domains and adjoint pairings.}\par
{\footnotesize \textit{Remaining.} Formalize the resolvent identity on the compatible closed-form domains.}\par
\smallskip
\noindent\textbf{\texttt{\small DERIV:\allowbreak{}YANGMILLS\_\allowbreak{}WEIGHTED\_\allowbreak{}CURVATURE:\allowbreak{}GST3}} \hfill {\footnotesize conditional}\par
{\small Ric+(2/hbar)Hess S>=kappa gives ||A\_\allowbreak{}S f||\textasciicircum{}2>=c*kappa<f,A\_\allowbreak{}Sf>, hence spectrum in \{0\} union [c*kappa,infinity) and the Poincare bound when the kernel is constants.}\par
{\footnotesize \srcpath{docs/derivations/yangmills-weighted-curvature.md} --- \#\# GST-3: curvature normalization and a spectral argument without a discrete excited state (lines 92-130)}\par
{\footnotesize \textit{Hypotheses.} Finite-dimensional smooth Riemannian configuration space or the explicitly specified finite matrix/oscillator model.; H=-c Delta+V, c=hbar\textasciicircum{}2/(2m)>0; psi=N exp(-S/hbar)>0 normalized; common operator/form domains as stated.; kappa>0; closed Bochner identity/inequality and simple constant kernel.}\par
{\footnotesize \textit{Remaining.} The bounded positive-operator square-to-gap implication is formalized in arbitrary complex Hilbert dimension. Construct the actual weighted operator, verify its quadratic inequality and manage the unbounded-domain/spectral realization required by the source.}\par
\smallskip
\noindent\textbf{\texttt{\small DERIV:\allowbreak{}YANGMILLS\_\allowbreak{}WEIGHTED\_\allowbreak{}CURVATURE:\allowbreak{}GST6}} \hfill {\footnotesize conditional}\par
{\small The closed-form min-max comparison yields E1(H)-E0(H)>=gamma-(mean R-rmin).}\par
{\footnotesize \srcpath{docs/derivations/yangmills-weighted-curvature.md} --- \#\# GST-6: the residual-mean certificate (lines 167-190)}\par
{\footnotesize \textit{Hypotheses.} Finite-dimensional smooth Riemannian configuration space or the explicitly specified finite matrix/oscillator model.; H=-c Delta+V, c=hbar\textasciicircum{}2/(2m)>0; psi=N exp(-S/hbar)>0 normalized; common operator/form domains as stated.; A\_\allowbreak{}S is nonnegative self-adjoint with simple normalized constant ground and first variational level >=gamma>0; R>=rmin, finite mean; 1 belongs to the sum-form domain; min-max applies.}\par
{\footnotesize \textit{Remaining.} Formalize the lower first-excited variational comparison and the trial-vector Rayleigh upper bound on common form domains.}\par
\smallskip
\endgroup
